\documentclass[11pt]{article}
\usepackage{amsmath,amsthm,amssymb}
\usepackage[margin=1in]{geometry}
\usepackage{hyperref}
\usepackage{booktabs}
\usepackage{enumitem}

\newtheorem{theorem}{Theorem}
\newtheorem{lemma}[theorem]{Lemma}
\newtheorem{proposition}[theorem]{Proposition}
\newtheorem{corollary}[theorem]{Corollary}
\newtheorem{conjecture}[theorem]{Conjecture}
\theoremstyle{definition}
\newtheorem{definition}[theorem]{Definition}
\newtheorem{remark}[theorem]{Remark}
\newtheorem{observation}[theorem]{Observation}
\newtheorem{question}[theorem]{Question}

\title{The induced-module image is not $B$:\\
  $L \subseteq \ker M$ via a one-line sign argument}
\author{Clio}
\date{14 May 2026 (afternoon)}

\begin{document}
\maketitle

\begin{abstract}
We retain notation: $\lambda \vdash n$ with $\ell = \ell(\lambda)$
rows; $H_q(S_n)$ acting on the Specht module $V_\lambda$ in the
Hoefsmit seminormal basis; $R'_j := (T_{j-1}+1)\cdots(T_1+1)$;
$M := R'_{\ell+1} R'_{\ell+2} \cdots R'_n$;
$B := M\,V_\lambda$; and the column-1-deletion shape
$\lambda^{(\ge 2)} := (\lambda_1-1, \lambda_2-1, \ldots)$.

The May-14 \textsc{Prove.md} conjectured that $B$ equals the image of
the canonical Frobenius extension
\[
   \psi : \mathrm{Ind}_{H_q(S_\ell) \times H_q(S_{n-\ell})}^{H_q(S_n)}
            \!\bigl( V_{(1^\ell)} \boxtimes V_{\lambda^{(\ge 2)}} \bigr)
        \;\longrightarrow\; V_\lambda
\]
of the LR-multiplicity-$1$ embedding. We refute this in the strongest
possible sense:

\textbf{Refutation.} For irreducible $V_\lambda$ at generic $q$,
$\mathrm{im}(\psi) = V_\lambda$ entirely (since the image is an
$H_q(S_n)$-submodule and $\psi \ne 0$). In every tested stable case,
$\dim \mathrm{im}(\psi) = \dim V_\lambda \gg \dim B$ (e.g.\ $64$ vs.\ $2$
for $\lambda = (3,2,1,1,1)$ at $n=8$).

\textbf{Stronger structural finding.} The LR-isotypic component
\[
   L \;:=\; V_{(1^\ell)} \boxtimes V_{\lambda^{(\ge 2)}}
   \;\hookrightarrow\; \mathrm{Res}_P V_\lambda
\]
sits in the \emph{kernel} of $M$, not in $B$:
$L \subseteq \ker M$. Indeed, more generally:

\begin{theorem}[Sign-kill lemma]\label{thm:sign-kill}
$E_1^-\!\mid_{V_\lambda} \;\subseteq\; \ker M$, where
$E_1^- = \ker(T_1 + 1)$ is the $(-1)$-eigenspace of $T_1$.
\end{theorem}

The proof is one line (each $R'_j$ ends with $(T_1+1)$ on the right;
hence $R'_n$ kills $E_1^-$ and so does $M$).

This is a genuine \emph{antipode} relationship: the subspace one would
naively map to $B$ via the leftmost-factor projector $M$ is precisely
the subspace $M$ annihilates. Computationally we observe further
$L \cap B = 0$ in every tested case, confirming that $L$ and $B$ are
disjoint dim-$f^{\lambda^{(\ge 2)}}$ subspaces of $V_\lambda$ ---
related by their dimension match (an LR coincidence) but otherwise
unrelated. The induced-module route (\textsc{May-13} resolution (c))
is therefore closed off; \emph{any} representation-theoretic
identification of $B$ with $V_{\lambda^{(\ge 2)}}$ must use machinery
other than the canonical Frobenius extension.
\end{abstract}

\section{Setup}\label{sec:setup}

We retain the conventions of the May-13 multi-corner paper and the
May-14 JM/non-stable note. Briefly:
\begin{itemize}[topsep=0.2em,itemsep=0.1em,leftmargin=2em]
\item $H_q(S_n)$ has generators $T_1, \ldots, T_{n-1}$ subject to
  $(T_i - q)(T_i + 1) = 0$ and the braid relations.
\item $V_\lambda$ is the Specht module in the Hoefsmit seminormal
  basis $\{v_T : T \in \mathrm{SYT}(\lambda)\}$.
\item For each $j \in \{1, \ldots, n-1\}$, $E_j^\pm \subseteq V_\lambda$
  is the $(\pm)$-eigenspace of $T_j$ (with eigenvalues $q$ and $-1$
  respectively).
\item $R'_j := (T_{j-1}+1)(T_{j-2}+1)\cdots(T_1+1)$ for $j \ge 2$,
  and $R'_1 := 1$.
\item For $\lambda$ with $\ell = \ell(\lambda)$ rows,
  $M := R'_{\ell+1} R'_{\ell+2} \cdots R'_n$, and
  $B := \mathrm{im}(M) = M \cdot V_\lambda$.
\item $\lambda^{(\ge 2)} := (\lambda_1 - 1, \lambda_2 - 1, \ldots,
  \lambda_\ell - 1)$ (drop zero parts), the column-$1$-deletion shape.
\item $P := H_q(S_\ell) \times H_q(S_{n-\ell})$ is the standard
  parabolic (acting on positions $\{1, \ldots, \ell\}$ and
  $\{\ell+1, \ldots, n\}$).
\end{itemize}

In the stable regime (May-13 closed-form theorem),
$\dim B = f^{\lambda^{(\ge 2)}}$.

\section{Refutation: the image of $\psi$ is all of $V_\lambda$}
\label{sec:refute-im}

The May-14 \textsc{Prove.md} proposed:
\begin{conjecture}[\textsc{May-14 Prove.md} target]\label{conj:prove}
In the stable regime, the canonical $H_q(S_n)$-equivariant morphism
\[
   \psi : \mathrm{Ind}_P^{H_q(S_n)}\!\bigl( V_{(1^\ell)}
        \boxtimes V_{\lambda^{(\ge 2)}} \bigr) \;\longrightarrow\; V_\lambda
\]
has image equal to $B$.
\end{conjecture}

\subsection{Why $\psi$ exists uniquely up to scalar}

By Frobenius reciprocity for finite-dimensional Hopf-like algebras
(in particular for $H_q(S_n)$, which is Frobenius and semisimple at
generic $q$),
\[
   \mathrm{Hom}_{H_q(S_n)}\!\bigl(\mathrm{Ind}_P^{H_q(S_n)}(V_{(1^\ell)} \boxtimes V_{\lambda^{(\ge 2)}}),\; V_\lambda \bigr)
   \;\cong\; \mathrm{Hom}_P\!\bigl( V_{(1^\ell)} \boxtimes V_{\lambda^{(\ge 2)}},\; \mathrm{Res}_P V_\lambda \bigr).
\]
The right-hand side has dimension equal to the multiplicity of
$V_{(1^\ell)} \boxtimes V_{\lambda^{(\ge 2)}}$ in $\mathrm{Res}_P V_\lambda$, which by
Littlewood--Richardson equals the LR coefficient
$c^\lambda_{(1^\ell),\,\lambda^{(\ge 2)}}$. This coefficient is~$1$
(the skew shape $\lambda / (1^\ell)$ is a vertical strip, with the
unique LR filling assigning content $\lambda^{(\ge 2)}$). Hence $\psi$
exists and is unique up to scalar.

\subsection{The image of $\psi$}

\begin{theorem}[Refutation]\label{thm:refute}
At generic $q$, $\mathrm{im}(\psi) = V_\lambda$. In particular, for
every $\lambda$ with $\dim V_\lambda > \dim B$ (which holds whenever
$\lambda$ is not an unusually small shape), $\mathrm{im}(\psi) \ne B$,
refuting Conjecture~\ref{conj:prove}.
\end{theorem}

\begin{proof}
$\mathrm{im}(\psi) \subseteq V_\lambda$ is an $H_q(S_n)$-submodule. At
generic $q$, $H_q(S_n)$ is semisimple and $V_\lambda$ is irreducible,
so the only submodules of $V_\lambda$ are $0$ and $V_\lambda$. Since
$\psi \ne 0$ (LR multiplicity~$1$), $\mathrm{im}(\psi)$ is a nonzero
submodule of $V_\lambda$, hence equal to $V_\lambda$. \qed
\end{proof}

\begin{remark}[Quantitative gap]
For $\lambda = (3,2,1,1,1)$ at $n = 8$ (a stable shape),
$\dim V_\lambda = 64$ and $\dim B = 2$. So
$\dim \mathrm{im}(\psi) = 64 \ne 2 = \dim B$. Similarly for every
larger stable case in the May-13 family table.
\end{remark}

The refutation is purely structural: the image of an equivariant map
out of an induced module is always an honest submodule of $V_\lambda$,
which by irreducibility forces it to be either zero or all of
$V_\lambda$. The dimension match $\dim B = \dim L = f^{\lambda^{(\ge 2)}}$
is real but does not lift to an iso of $\psi$-images.

\section{The sign-kill lemma}\label{sec:sign-kill}

\begin{theorem}[Sign-kill]\label{thm:sign-kill-body}
For any $\lambda \vdash n$ with $\ell = \ell(\lambda)$,
\[
   E_1^-\!\mid_{V_\lambda} \;\subseteq\; \ker M.
\]
That is, every vector $v \in V_\lambda$ with $T_1 v = -v$ satisfies
$M v = 0$.
\end{theorem}

\begin{proof}
The leftmost factor of $R'_j$ in its standard expansion is $(T_{j-1}+1)$;
the rightmost factor is $(T_1 + 1)$. Hence for $j \ge 2$,
\[
   R'_j \;=\; X_j \cdot (T_1 + 1), \qquad
   X_j := (T_{j-1}+1)(T_{j-2}+1)\cdots(T_2+1) \;\in\; H_q(S_n).
\]
For $v \in E_1^-$ (so $T_1 v = -v$), $(T_1 + 1) v = 0$, and therefore
$R'_j v = X_j \cdot 0 = 0$ for every $j \ge 2$. In particular,
$R'_n v = 0$, so
\[
   M v \;=\; R'_{\ell+1} R'_{\ell+2} \cdots R'_{n-1} \cdot R'_n v
        \;=\; R'_{\ell+1} \cdots R'_{n-1} \cdot 0 \;=\; 0. \qed
\]
\end{proof}

\subsection{Corollaries}

\begin{corollary}[$L \subseteq \ker M$]\label{cor:L-in-ker}
The LR-isotypic component
$L := V_{(1^\ell)} \boxtimes V_{\lambda^{(\ge 2)}} \subseteq
\mathrm{Res}_P V_\lambda$ satisfies $L \subseteq \ker M$.
\end{corollary}

\begin{proof}
$V_{(1^\ell)}$ is the sign representation of $H_q(S_\ell)$, on which
$T_i$ acts as $-1$ for every $i \in \{1, \ldots, \ell - 1\}$. Hence
the LR-isotypic component, viewed inside $V_\lambda$, satisfies
$T_i v = -v$ for all $i \in \{1, \ldots, \ell - 1\}$ and $v \in L$.
In particular, $T_1 v = -v$, so $L \subseteq E_1^-$. By
Theorem~\ref{thm:sign-kill-body}, $L \subseteq \ker M$. \qed
\end{proof}

\begin{corollary}[Combinatorial description of $L$ in Hoefsmit basis]
\label{cor:L-Hoefsmit}
In the Hoefsmit seminormal basis,
\[
   L \;=\; \mathrm{span}\!\bigl\{ v_T : T \in \mathrm{SYT}(\lambda),\
      T(i, 1) = i \text{ for } i = 1, \ldots, \ell \bigr\}.
\]
That is, $L$ is the span of $v_T$ for SYTs $T$ whose column~$1$ is
filled by $1, 2, \ldots, \ell$ in order; equivalently, $T|_{[\ell]}$
is the unique SYT of shape $(1^\ell)$. In particular,
$\dim L = f^{\lambda / (1^\ell)} = f^{\lambda^{(\ge 2)}}$.
\end{corollary}

\begin{proof}
By Hoefsmit, the parabolic isotypic
$S_\mu := \mathrm{span}\{v_T : T|_{[\ell]} \text{ has shape } \mu\}$
of $V_\lambda \!\downarrow_P$ is exactly $V_\mu \boxtimes V^{\lambda/\mu}$,
where $V^{\lambda/\mu}$ is the skew Specht module. For $\mu = (1^\ell)$,
$V^{\lambda/(1^\ell)} \cong V_{\lambda^{(\ge 2)}}$ as $H_q(S_{n-\ell})$-modules
(by the LR rule, the skew shape $\lambda/(1^\ell)$ has $V_{\lambda^{(\ge 2)}}$
as its unique decomposition summand, with multiplicity~$1$). Hence
$S_{(1^\ell)} \cong V_{(1^\ell)} \boxtimes V_{\lambda^{(\ge 2)}} = L$ as
$P$-modules, and the dimension count is the SYT count of the skew
shape $\lambda / (1^\ell)$, which by translation equals
$f^{\lambda^{(\ge 2)}}$. \qed
\end{proof}

\subsection{$L$ and $B$ are disjoint in every tested case}

\begin{observation}\label{obs:disjoint}
In every stable case computed (Section~\ref{sec:data}),
\[
   L \cap B \;=\; 0.
\]
That is, $L$ and $B$ are two distinct dim-$f^{\lambda^{(\ge 2)}}$
subspaces of $V_\lambda$ (with $L \subseteq \ker M$ by
Corollary~\ref{cor:L-in-ker}, and $B = \mathrm{im}(M)$ by definition),
and they meet only at zero. The structural reason for this disjointness
is presently unproven; the most natural attempt
($V_\lambda = \ker M \oplus \mathrm{im}(M)$, which would force
$L \cap B \subseteq \ker M \cap \mathrm{im}(M) = 0$) fails in general
because the Fitting decomposition of $M$ is non-trivial: in some
stable cases $M^2 = 0$ (e.g.\ $\lambda = (3,2,1,1,1)$, where $B
\subseteq \ker M$ as well), and in others $\mathrm{rk}(M^2) =
\mathrm{rk}(M)$ but $M$ is not idempotent up to scalar. The
disjointness $L \cap B = 0$ holds in both regimes.
\end{observation}

\section{The strict-superset structure of $\ker M$}
\label{sec:ker-structure}

Theorem~\ref{thm:sign-kill-body} gives a sharp \emph{lower bound}:
$\ker M \supseteq E_1^-$. The kernel is in fact strictly larger;
the gap is what makes $\dim B$ much smaller than $\dim E_1^+$.

\begin{remark}[$E_i^- \not\subseteq \ker M$ for $i \ge 2$]
\label{rem:Ei-not-killed}
A computational check (Section~\ref{sec:data}) shows that
$E_i^- \not\subseteq \ker M$ for any $i \ge 2$: in fact
$\dim M(E_i^-) = \dim B$ for all $i \ge 2$ in the tested stable
cases. So the sign-kill phenomenon is genuinely tied to position~$1$
and does not extend to other adjacent transpositions.
\end{remark}

\begin{remark}[The kernel via parabolic isotypics]
\label{rem:para-kernel}
Decomposing $V_\lambda \!\downarrow_P = \bigoplus_{\mu \vdash \ell,\,
\mu \subseteq \lambda} S_\mu$ with $S_\mu = V_\mu \boxtimes V^{\lambda/\mu}$,
computational scan reveals:
\[
   M(S_\mu) \;=\; 0 \quad\text{for some structurally identifiable family of }\mu.
\]
Specifically, $S_{(1^\ell)} = L \subseteq \ker M$ (this is
Corollary~\ref{cor:L-in-ker}), and the data shows that
$S_{(2, 1^{\ell-2})} \subseteq \ker M$ as well (in every tested case),
along with several further $\mu$'s in larger shapes. The exact
combinatorial criterion is open; the data is recorded in
Table~\ref{tab:para-data}.
\end{remark}

\begin{table}[t]
\centering
\renewcommand{\arraystretch}{1.15}
\begin{tabular}{l|l|l}
\toprule
$\lambda$ & killed $\mu$'s & contributing $\mu$'s \\
\midrule
$(3, 2, 1)$         & $(1^3)$                                & $(2, 1),\ (3)$ \\
$(3, 2, 1, 1)$      & $(1^4),\ (2, 1, 1)$                    & $(2, 2),\ (3, 1)$ \\
$(3, 2, 1, 1, 1)$   & $(1^5),\ (2, 1^3)$                     & $(2, 2, 1),\ (3, 1, 1),\ (3, 2)$ \\
$(3, 3, 1, 1)$      & $(1^4),\ (2, 1, 1)$                    & $(2, 2),\ (3, 1)$ \\
$(4, 2, 1^4)$       & $(1^6),\ (2, 1^4),\ (2, 2, 1, 1),\ (3, 1, 1, 1)$ & $(3, 2, 1),\ (4, 1, 1),\ (4, 2)$ \\
$(4, 3, 1^3)$       & $(1^5),\ (2, 1^3)$                     & $(2, 2, 1),\ (3, 1, 1),\ (3, 2),\ (4, 1)$ \\
\bottomrule
\end{tabular}
\caption{The parabolic decomposition $V_\lambda \!\downarrow_P = \bigoplus_\mu S_\mu$:
which $\mu$ have $M(S_\mu) = 0$ (\emph{killed}) and which contribute
nontrivially to $B = \sum_\mu M(S_\mu)$ (\emph{contributing}). Computed
mod $P = 100\,003$ at $q = 11$ via
\texttt{probe\_parabolic.py}.}
\label{tab:para-data}
\end{table}

\begin{conjecture}[Parabolic kill criterion, partial]\label{conj:para}
There is a combinatorial criterion on $\mu$ (depending on $\lambda$)
such that $S_\mu \subseteq \ker M$ for every $\mu$ failing the
criterion. The criterion includes at minimum
$\mu \in \{(1^\ell),\ (2, 1^{\ell-2})\}$, since these always satisfy
$T_1 v = -v$ (or have a related reduction to the same). A precise
statement is open.
\end{conjecture}

\section{Implications for the iso programme}\label{sec:implications}

\begin{enumerate}[leftmargin=2em,topsep=0.3em,itemsep=0.3em]
\item \textbf{The induced-module route is closed.} The May-13
``resolution (c)'' suggested that $B$ might be the image of an
equivariant morphism out of an induced module containing
$V_{(1^\ell)} \boxtimes V_{\lambda^{(\ge 2)}}$. Theorem~\ref{thm:refute}
shows that the obvious such morphism (the canonical Frobenius
extension $\psi$) has image equal to all of $V_\lambda$, not~$B$.
Combined with the May-13 refutation of standard parabolic embeddings
and the May-14 refutation of seminormal-subset spans, this exhausts
the natural representation-theoretic candidates for an iso
$B \cong V_{\lambda^{(\ge 2)}}$ via parabolic-style induction or
restriction.

\item \textbf{The antipode picture.} The LR-isotypic $L$ is the
``natural input'' to $\psi$ (it is the embedded copy of
$V_{(1^\ell)} \boxtimes V_{\lambda^{(\ge 2)}}$ in $V_\lambda \!\downarrow_P$).
Yet $L \subseteq \ker M$, so applying the leftmost-factor projector
$M$ to $L$ gives $0$. Meanwhile $B = \mathrm{im}(M)$ is a different
subspace of the same dimension. This is the strongest possible
mismatch: $L$ is exactly what $M$ \emph{annihilates}, while $B$ is
exactly what $M$ \emph{produces}. The dimension match is therefore
a coincidence of LR-counting and $f^{\lambda^{(\ge 2)}}$-counting,
with no direct map between $L$ and $B$ via the standard induced-module
machinery.

\item \textbf{What might still work.} Three remaining routes are not
obviously closed:
\begin{enumerate}[label=(\roman*),topsep=0.2em,itemsep=0.1em]
\item \emph{The dual Frobenius map.} The companion morphism
$\phi : V_\lambda \to \mathrm{Ind}_P^{H_q(S_n)}(V_{(1^\ell)}
\boxtimes V_{\lambda^{(\ge 2)}})$ exists by the other Frobenius
adjunction, with $\dim \mathrm{Hom}_G(V_\lambda, \mathrm{Ind}) = c^\lambda_{(1^\ell),\,\lambda^{(\ge 2)}} = 1$.
By irreducibility, $\phi$ is injective. Whether $\phi^{-1}$ of some
natural subspace of $\mathrm{Ind}$ equals $B$ is open.

\item \emph{Cellular / KL bases.} The Hecke algebra $H_q(S_n)$ has a
cellular basis (Kazhdan--Lusztig), and the seminormal Hoefsmit basis
is one of several natural choices. The chains of the May-13 paper
might lift to a different cellular description of $B$.

\item \emph{Subquotient construction.} $B$ might admit a natural
description as a subquotient of $V_\lambda$ under a non-standard
filtration, and the graded piece carrying $V_{\lambda^{(\ge 2)}}$ might
be isomorphic to (a twist of) $V_{\lambda^{(\ge 2)}}$ as an
$H_q(S_{n-\ell})$-module under \emph{some} embedding.
\end{enumerate}
\end{enumerate}

\section{Computational verification}\label{sec:data}

All computations were performed mod $P = 100\,003$ at $q = 11$ using
the seminormal-matrix infrastructure of the May-13 papers. Mod-$P$
ranks coincide with $\mathbb{Q}(q)$-ranks at any non-special $q$;
$q = 11$ is non-special for $|\lambda| \le 11$.

\subsection{Verification of $L \subseteq \ker M$ and $L \cap B = 0$}

\begin{table}[h]
\centering
\renewcommand{\arraystretch}{1.15}
\begin{tabular}{l|c|c|c|c|c|c}
\toprule
$\lambda$ & $n$ & $\dim V$ & $\dim L$ & $\dim B$ & $\dim M(L)$ & $\dim(L + B)$ \\
\midrule
$(3, 2, 1)$        & 6  & 16  & 2 & 2 & 0 & 4 \\
$(3, 2, 1, 1)$     & 7  & 35  & 2 & 2 & 0 & 4 \\
$(3, 2, 1, 1, 1)$  & 8  & 64  & 2 & 2 & 0 & 4 \\
$(3, 3, 1, 1)$     & 8  & 56  & 2 & 2 & 0 & 4 \\
$(4, 2, 1^4)$      & 10 & 350 & 3 & 3 & 0 & 6 \\
$(4, 3, 1^3)$      & 10 & 525 & 5 & 5 & 0 & 10 \\
\bottomrule
\end{tabular}
\caption{$\dim M(L) = 0$ in every case (confirming
Corollary~\ref{cor:L-in-ker}), and $\dim(L + B) = \dim L + \dim B$
in every case (confirming $L \cap B = 0$,
Observation~\ref{obs:disjoint}).}
\label{tab:LB-data}
\end{table}

\subsection{Verification of $E_1^- \subseteq \ker M$}

\begin{table}[h]
\centering
\renewcommand{\arraystretch}{1.15}
\begin{tabular}{l|c|c|c|c}
\toprule
$\lambda$ & $\dim V$ & $\dim B$ & $\dim E_1^-$ & $\dim M(E_1^-)$ \\
\midrule
$(3, 2, 1)$        & 16 & 2 & 8  & 0 \\
$(3, 2, 1, 1)$     & 35 & 2 & 20 & 0 \\
$(3, 2, 1, 1, 1)$  & 64 & 2 & 40 & 0 \\
$(3, 3, 1, 1)$     & 56 & 2 & 30 & 0 \\
\bottomrule
\end{tabular}
\caption{$M(E_1^-) = 0$ in every case (verifying
Theorem~\ref{thm:sign-kill-body} computationally).}
\end{table}

\subsection{Failure of $E_i^-$ for $i \ge 2$}

For the same shapes, $\dim M(E_i^-) = \dim B$ for every $i \in \{2, 3,
\ldots, n-1\}$ --- i.e., the operator $M$ surjects from $E_i^-$ onto
$B$ for $i \ge 2$, so $E_i^-$ is \emph{not} contained in $\ker M$.

\subsection{Scripts}

The computations were performed by:
\begin{itemize}[leftmargin=2em,topsep=0.2em,itemsep=0.1em]
\item \texttt{/home/clio/projects/scratch/2026-05-14-induced-image/compare\_L\_B.py}
  --- computes $L$, $B$, $M(L)$, $L \cap B$ for each test case.
\item \texttt{/home/clio/projects/scratch/2026-05-14-induced-image/probe\_parabolic.py}
  --- decomposes $V_\lambda \!\downarrow_P$ into $S_\mu$'s and tests
  $\dim M(S_\mu)$ for each $\mu$.
\item \texttt{/home/clio/projects/scratch/2026-05-14-induced-image/probe\_E\_minus.py}
  --- builds $E_i^- = \ker(T_i + 1)$ and tests $\dim M(E_i^-)$ for each $i$.
\item \texttt{/home/clio/projects/scratch/2026-05-14-induced-image/probe\_M\_squared.py}
  --- tests Fitting structure of $M$ ($\mathrm{rk}\,M^2$ vs.\
  $\mathrm{rk}\,M$).
\end{itemize}

\section{Honest status}\label{sec:status}

This note's main content is a refutation of the \textsc{May-14
Prove.md} target (Conjecture~\ref{conj:prove}) at the strongest
possible level: not only is $\mathrm{im}(\psi) \ne B$, but
$\mathrm{im}(\psi) = V_\lambda$ entirely, and the would-be input $L$
is precisely killed by $M$. Combined with the May-13 and May-14
refutations of standard parabolic embeddings and seminormal-subset
spans, the natural ``inductive / restricted Schur--Weyl''
representation-theoretic identification of $B$ with $V_{\lambda^{(\ge
2)}}$ is now ruled out via every routine candidate.

The positive content is the one-line sign-kill lemma
(Theorem~\ref{thm:sign-kill-body}), which makes the structural reason
for $L \subseteq \ker M$ transparent and explains the parabolic
kill of $S_{(1^\ell)}$ uniformly. The further parabolic kill data
(Table~\ref{tab:para-data}) suggests a richer combinatorial criterion
worth pursuing.

The remaining open route is the dual-Frobenius map $\phi: V_\lambda
\to \mathrm{Ind}_P^G(V_{(1^\ell)} \boxtimes V_{\lambda^{(\ge 2)}})$,
which is injective by irreducibility; whether some natural pre-image
under $\phi$ equals $B$ is the leading remaining candidate, untested
here.

\appendix
\section{Why the standard symmetrizer-style identities don't help}\label{app:why-not}

For a Hecke algebra $H_q(S_k)$ with generators $T_1, \ldots, T_{k-1}$,
the \emph{trivial idempotent} (up to a scalar) admits the closed-form
factorisation
\[
   h_k^+ \;=\; (T_{k-1}+1)(T_{k-2} T_{k-1} + T_{k-1} + T_{k-2} + 1) \cdots
\]
or one of several Matsumoto-style products over the longest element.
The element $R'_n = (T_{n-1}+1)\cdots(T_1+1)$ used here is \emph{not}
the trivial idempotent; it is one factor of length $n - 1$ (rather
than $\binom{n}{2}$ as required), and its image is genuinely smaller
than the trivial isotypic of $V_\lambda \!\downarrow_{H_q(S_n)}$.

Concretely, $R'_n$ does not even square to a scalar multiple of
itself in general; the May-14 \texttt{probe\_M\_squared.py} data
shows $M^2 = 0$ for $\lambda = (3,2,1,1,1)$ (so $B \subseteq \ker M$
in that case) and $\mathrm{rk}\,M^2 = \mathrm{rk}\,M$ but $M^2 \ne c
M$ for $\lambda = (3,2,1)$. The precise algebraic structure of $M$
inside $H_q(S_n)$ is shape-dependent.

\end{document}
